% =============================================================================
% Section 4 — Cryptographic Layer
% =============================================================================
\section{Cryptographic Layer}
\label{sec:crypto}

\pebble{} uses Ed25519~\cite{bernstein2012ed25519,rfc8032} for all digital
signatures. Ed25519 offers 128-bit security, fast verification, small 64-byte
signatures, and no per-signature randomness requirement, making it well-suited
for a high-throughput distributed system.

\subsection{Node Identity and Key Derivation}

Each node possesses a unique cryptographic identity generated once and persisted
to \texttt{\{dataDir\}/identity.json}.

\begin{definition}[Node identity]
A node's identity is a triple $\mathcal{I} = (\mathrm{nodeId},\, pk,\, sk)$
where:
\begin{itemize}[leftmargin=1.5em]
  \item $(pk, sk)$ is an Ed25519 key pair. $pk \in \{0,1\}^{256}$,
        $sk \in \{0,1\}^{256}$ (the seed).
  \item $\mathrm{nodeId} = \mathrm{base64url}\bigl(\mathrm{SHA\text{-}256}(pk)\bigr)$
        is the globally unique node identifier, derived deterministically from the
        public key.
\end{itemize}
\end{definition}

The identity file is created on first start and loaded on subsequent starts,
ensuring the same \texttt{nodeId} across restarts. The private key seed is
stored in hex encoding and never transmitted over any network interface.

\subsection{Canonical JSON Serialisation}

To produce a deterministic byte sequence for signing, \pebble{} implements a
canonical JSON encoder inspired by RFC~8785~\cite{rfc8785}.

\begin{definition}[Canonical encoding]
\label{def:canonical}
$\mathrm{canonicalEncode}(v)$ serialises a JSON-compatible value $v$ to a
UTF-8 byte string such that:
\begin{enumerate}[leftmargin=1.5em]
  \item Object keys are sorted lexicographically by their UTF-16 code units
        (matching the ECMAScript \texttt{localeCompare} order).
  \item No insignificant whitespace is emitted.
  \item Numbers are represented in their shortest round-trip form.
  \item String escaping follows JSON specification.
\end{enumerate}
\end{definition}

The canonical encoder is the only source of truth for the byte sequence
presented to the signing primitive. Any serialisation discrepancy would cause
valid signatures to appear invalid; therefore the encoder is tested with a
deterministic golden-value test suite.

\subsection{Mutation Signing and Verification}

\begin{definition}[Signable mutation fields]
The \emph{signable fields} of a mutation $m$ are all fields except $\msig$
itself:
\[
  m_{\text{sign}} =
    \langle m.\mnid,\; m.\mseq,\; m.\mts,\; m.\mop,\; m.\mkey,\; m.\mval \rangle
\]
\end{definition}

\begin{protocol}[Mutation signing]
To create a signed mutation at node $n$ with identity $\mathcal{I}_n$:
\begin{enumerate}[leftmargin=1.5em]
  \item Assemble $m_{\text{sign}}$ with all fields populated.
  \item Compute $b = \mathrm{canonicalEncode}(m_{\text{sign}})$.
  \item Compute $\msig = \mathrm{Ed25519Sign}(sk_n,\, b)$.
  \item Return $m = m_{\text{sign}} \cup \{\msig\}$.
\end{enumerate}
\end{protocol}

\begin{protocol}[Mutation verification]
\label{proto:verify-mut}
To verify a mutation $m$ received from a peer, given the originating node's
public key $pk_{\mnid}$ obtained from the cluster configuration:
\begin{enumerate}[leftmargin=1.5em]
  \item Extract $m_{\text{sign}} = m \setminus \{\msig\}$.
  \item Compute $b = \mathrm{canonicalEncode}(m_{\text{sign}})$.
  \item Return $\mathrm{Ed25519Verify}(pk_{\mnid},\, b,\, m.\msig)$.
\end{enumerate}
If the originating \texttt{nodeId} is not present in the cluster configuration,
verification fails with an \texttt{UnknownNodeError} before the signature is
even checked.
\end{protocol}

\begin{remark}[Startup recovery bypass]
During crash recovery (\cref{sec:persistence}), mutations are replayed from the
local on-disk log with signature verification \emph{skipped}. This is safe
because:
(1)~the log is written by the local node and protected by OS file permissions;
(2)~any mutation in the log was either created locally (and therefore valid) or
received from a peer and verified at application time before being persisted.
The \texttt{skipSignatureVerification} flag is never set for mutations received
over the network.
\end{remark}

\subsection{Moderator Credential System}

External client access is controlled through a lightweight certificate-authority
(CA) model. The \emph{moderator} is an offline entity that holds the cluster's
root signing key.

\begin{definition}[Client credential]
\label{def:credential}
A \emph{client credential} is a 6-tuple signed by the moderator:
\[
  C = \langle \mathit{clientId},\; pk_c,\; \mathit{perm},\;
              t_{\text{issued}},\; t_{\text{expires}},\; \msig_{\text{mod}} \rangle
\]
where:
\begin{itemize}[leftmargin=1.5em]
  \item $\mathit{clientId}$ is a human-readable client identifier.
  \item $pk_c \in \{0,1\}^{256}$ is the client's Ed25519 public key.
  \item $\mathit{perm} \in \{\texttt{read-only},\; \texttt{read-write}\}$ is
        the access level granted.
  \item $t_{\text{issued}},\; t_{\text{expires}} \in \mathbb{Z}$ are Unix
        millisecond timestamps ($t_{\text{expires}} = \bot$ for no expiry).
  \item $\msig_{\text{mod}} = \mathrm{Ed25519Sign}\bigl(sk_{\text{mod}},\;
        \mathrm{canonicalEncode}(\langle\mathit{clientId}, pk_c, \mathit{perm},
        t_{\text{issued}}, t_{\text{expires}}\rangle)\bigr)$.
\end{itemize}
\end{definition}

Credential issuance is performed offline using \pkg{moderator} and does not
require modifying \texttt{cluster.json} or restarting running nodes. The
moderator's public key $pk_{\text{mod}}$ is stored in the cluster configuration
and used by every node to verify incoming credentials.

\begin{remark}[Security properties]
The credential system provides:
\begin{itemize}[leftmargin=1.5em]
  \item \textbf{Unforgeability}: only the holder of $sk_{\text{mod}}$ can issue
        valid credentials.
  \item \textbf{Expiry enforcement}: nodes reject credentials where
        $\mathrm{now}() > t_{\text{expires}}$.
  \item \textbf{Key binding}: during the AUTH handshake, the client proves
        possession of $sk_c$ by signing a fresh nonce, preventing credential
        replay from a different identity.
\end{itemize}
\end{remark}
